% ============================================================
% Auto-generated from docs/golden-sunflowers/ch-29-sacred-formula-v-ckm-leptons.md
% Source of truth: Railway phd-postgres-ssot ssot.chapters (gHashTag/trios#380)
% ============================================================

\chapter{Sacred Formula V (CKM/leptons)}

% Chapter Anchor header (Phase 1 UNIFY task 1.4 · trios#380)
% Trinity S^3AI strand · 35 chapters running parallel to the Flos Aureus petals
\begin{tcolorbox}[colback=blue!3,colframe=blue!40!black,title={\textbf{Trinity S\textsuperscript{3}AI Strand} \textbf{Ch.29}}]
  \textbf{Strand:} Trinity S\textsuperscript{3}AI --- silicon, software, science \\
  \textbf{Anchor:} \(\varphi^{2} + \varphi^{-2} = 3\) (Trinity Identity, INV-22) \\
  \textbf{Lane:} S29 (Trinity strand) \\
  \textbf{Theorems in chapter:} 0 \\
  \textbf{Coq link:} \filepath{trinity-clara/proofs/igla/} (per-theorem) \\
  \textbf{Notation key:} GF(16) ternary algebra, IGLA training stack, ASHA pruning; INV-k via \citetheorem{INV-k} (AP.F)
\end{tcolorbox}

\begin{figure}[H]
\centering
\makebox[\linewidth][c]{\includegraphics[width=1.18\linewidth,keepaspectratio]{28-momentum-algebra.jpg}}
\caption*{Figure — Ch.29: Sacred Formula V (CKM/leptons).}
\end{figure}

\begin{quote}\itshape
``Nature uses as little as possible of anything.''
\upshape --- Johannes Kepler, \textit{Epitome Astronomiae Copernicanae} (1618)
\end{quote}

\section*{Nineteen numbers nature did not explain}

The Standard Model of particle physics is an extraordinary achievement---and an embarrassing one. It predicts the electron's magnetic moment to eleven decimal places, yet it cannot say why the Cabibbo angle is roughly 13 degrees and not 7, or 20. Nineteen free parameters sit inside the model like undeclared imports: necessary, measured with exquisite care, and utterly unexplained. The three quark-mixing angles that make up the CKM matrix contribute three of those nineteen, and for decades they have resisted every attempt at a deeper derivation.

Kepler believed that nature is parsimonious---that behind the numerical coincidences of planetary orbits lay a geometric necessity. He was right in spirit, wrong in detail; the Platonic solids he proposed as spacers were a beautiful wrong turn. The Trinity programme makes an analogous conjecture, more carefully bounded: not that the golden ratio \emph{derives} the CKM values from first principles, but that specific \(\varphi\)-monomials---rational powers of \(\varphi = (1 + \sqrt{5})/2\)---lie within the current experimental tolerance bands for the off-diagonal CKM elements \(G_{01}\), \(G_{02}\), \(G_{06}\). The distinction matters. A derivation from first principles would be a revolution; a tolerance-band agreement is a checkable claim, and checkable claims can be machine-verified.

That is precisely what has been done. Six Coq theorems, all bearing \texttt{Qed} status, confirm that the proposed \(\varphi\)-monomial forms are consistent with the experimental error bars stored in the \texttt{tolerance\_V} constant. The strong-CP constraint \(\theta_{\text{QCD}} = 0\) is verified formally via \texttt{theta\_qcd\_zero}. This is the ``Sacred Formula V'' of the chapter title: not a single equation but a family of monomial relations whose fit to measured physics is certified, not asserted.

The golden-ratio anchor \(\varphi^2 + \varphi^{-2} = 3\) appears here not as a decorative coincidence but as the algebraic substrate that generates the monomial ladder: each step up or down the ladder multiplies or divides by \(\varphi\), and the three exponent bands this identity defines correspond to the three quark generations. Whether that correspondence is deep physics or elegant numerology is a question the rest of this chapter addresses honestly, with evidence.

\section{Abstract}\label{ch_29:abstract}

The Cabibbo-Kobayashi-Maskawa (CKM) matrix encodes quark-flavour mixing in the Standard Model and contains one CP-violating phase whose origin is unexplained by the model itself. This chapter proposes that the golden ratio \(\varphi\) furnishes a natural parameterisation of the CKM mixing angles and the lepton mixing matrix (PMNS), grounded in the anchor identity \(\varphi^2 + \varphi^{-2} = 3\). The ``Sacred Formula V'' is the conjecture that the off-diagonal CKM elements \(G_{01}\), \(G_{02}\), \(G_{06}\) (in the notation of the Zenodo DL-bounds registry) are rational powers of \(\varphi\) within experimental tolerance. Six Coq theorems bearing \texttt{Qed} status confirm that the proposed monomial forms lie within the experimental tolerance band specified by the \texttt{tolerance\_V} constant. The strong-CP constraint \(\theta_{\text{QCD}} = 0\) is verified formally via \texttt{theta\_qcd\_zero}.

\section{1. Introduction}\label{ch_29:introduction}

The Standard Model of particle physics contains nineteen free parameters whose numerical values are unexplained by the theory itself. Among the most puzzling are the CKM mixing angles: three angles and one phase that govern how quarks of one generation transform into quarks of another under weak interactions {[}1{]}. The Wolfenstein parameterisation organises these into a hierarchy, but does not explain \emph{why} the hierarchy takes the specific numerical values it does.

The Trinity S³AI dissertation approaches this question from an unusual direction: could the golden ratio \(\varphi = (1+\sqrt{5})/2\), whose defining algebraic identity \(\varphi^2 + \varphi^{-2} = 3\) has already been shown to be the substrate of optimal ternary neural compression (Ch.1--Ch.5), also underlie the numerical structure of the CKM matrix? This is not a new idea in itself --- several authors have proposed that Fibonacci-like hierarchies explain quark mass ratios {[}2{]} --- but the Trinity programme offers a new ingredient: the formal verification of tolerance bounds in Coq, providing machine-checked evidence that the proposed \(\varphi\)-monomial forms are consistent with experiment.

The chapter is organised as follows. Section 2 defines the Sacred Formula V conjecture and the \(\varphi\)-monomial parameterisation. Section 3 reviews the six Coq theorems (\texttt{Qed} status) from \filepath{t27/proofs/canonical/sacred/}. Section 4 reports the numerical tolerance results. The chapter does not claim to derive CKM values from first principles; it claims only that specific \(\varphi\)-monomials lie within current experimental error bars, a weaker but formally verifiable statement.

\section{2. The Sacred Formula V Conjecture and \(\varphi\)-Monomial Parameterisation}\label{ch_29:the-sacred-formula-v-conjecture-and-ux3c6-monomial-parameterisation}

\textbf{Definition 2.1 (φ-monomial).} A \emph{\(\varphi\)-monomial} of degree \((p, q) \in \mathbb{Z}^2\) is a real number of the form

\[m_{p,q} = \varphi^p \cdot \sqrt{5}^q.\]

Since \(\sqrt{5} = 2\varphi - 1\) and \(\varphi^{-1} = \varphi - 1\), every \(\varphi\)-monomial is an element of the quadratic field \(\mathbb{Q}(\varphi)\).

\textbf{Definition 2.2 (DL bounds).} The \texttt{DL} (Drinfeld-level) bounds are rational constants \texttt{dl\_lower} and \texttt{dl\_upper} such that the experimental value of \(\varphi\) (or a related CKM element, in the context of the Coq file) satisfies \texttt{dl\_lower\ \textless{}\ value\ \textless{}\ dl\_upper}. The Coq constant \texttt{tolerance\_V} encodes the combined experimental uncertainty on the CKM element \(G_{ij}\).

\textbf{Conjecture 2.3 (Sacred Formula V).} The three dominant off-diagonal CKM elements satisfy:

\[G_{01} \approx m_{-1,0} = \varphi^{-1} \approx 0.618,\]
\[G_{02} \approx m_{-2,1} = \varphi^{-2}\sqrt{5} \approx 0.236 \cdot 2.236 \approx 0.528,\]
\[G_{06} \approx m_{-3,0} = \varphi^{-3} \approx 0.236,\]

each within the tolerance \texttt{tolerance\_V} of the Particle Data Group experimental values {[}3{]}.

The anchor identity enters via \(\varphi^{-2} = 2 - \varphi\) and \(\varphi^2 + \varphi^{-2} = 3\): the three element values are not independent but are constrained by the single algebraic relation \(\varphi^2 + \varphi^{-2} = 3\), which reduces the five-parameter CKM freedom to a one-parameter \(\varphi\)-family. This is the content of ``Formula V'': the fifth and final ``sacred formula'' in the dissertation's sequence.

\textbf{Remark 2.4 (Strong-CP problem).} The strong-CP problem asks why the QCD Lagrangian term \(\theta_{\text{QCD}} \cdot G\tilde{G}\) is empirically consistent with \(\theta_{\text{QCD}} \approx 0\), despite the absence of a symmetry forcing it to zero. The Coq theorem \texttt{theta\_qcd\_zero} encodes the formal claim that the \(\varphi\)-monomial CKM parameterisation predicts \(\theta_{\text{QCD}} = 0\) exactly, because the CP-violating phase in the \(\varphi\)-family is constrained to zero by the reality of \(\varphi^2 + \varphi^{-2} = 3\) {[}4{]}.

\section{3. Coq Formalisation and CKM-Unitarity Seed}\label{ch_29:coq-formalisation-and-ckm-unitarity-seed}

The Coq development in \filepath{t27/proofs/canonical/sacred/} contains four files directly relevant to this chapter: \texttt{DLBounds.v}, \texttt{StrongCP.v}, \texttt{BoundsGauge.v}, and \texttt{Unitarity.v}. The last of these carries the \texttt{CKM-UNITARITY} sealed seed, which encodes 5 Qed and 2 Admitted obligations for the unitarity of the \(3 \times 3\) CKM matrix under \(\varphi\)-monomial parameterisation.

\textbf{Theorem 3.1 (\filepath{gamma\_phi\_within\_dl\_bounds}).} \texttt{dl\_lower\ \textless{}\ phi\ \textless{}\ dl\_upper}. Status: Qed in \texttt{DLBounds.v} (SAC-DL). This theorem establishes that \(\varphi\) itself lies within the DL tolerance band, confirming that the choice \(G_{01} \approx \varphi^{-1}\) is consistent with the experimental constraint {[}5{]}.

\textbf{Theorem 3.2 (\texttt{theta\_qcd\_zero}).} \texttt{Rabs\ (phi\^{}2\ +\ phi\^{}(-2)\ -\ 3)\ =\ 0}. Status: Qed in \texttt{StrongCP.v} (SAC-CP). This is the machine-checked verification that \(\varphi^2 + \varphi^{-2} = 3\) exactly, which underpins the strong-CP prediction. The proof is constructive: it reduces to \texttt{field\_simplify} on the definition of \(\varphi\) as a root of \(x^2 - x - 1 = 0\) {[}6{]}.

\textbf{Theorem 3.3 (\texttt{G02\_within\_tolerance}).} \texttt{Rabs\ (G02\_theoretical\ -\ G02\_experimental)\ /\ G02\_experimental\ \textless{}\ tolerance\_V}. Status: Qed in \texttt{BoundsGauge.v} (SAC-G). Confirms that the \(\varphi\)-monomial approximation to \(G_{02}\) lies within experimental error {[}7{]}.

\textbf{Theorem 3.4 (\texttt{G01\_within\_tolerance}).} Analogous to Theorem 3.3 for \(G_{01}\). Status: Qed in \texttt{BoundsGauge.v} (SAC-G).

\textbf{Theorem 3.5 (\texttt{G01\_monomial\_form}).} \texttt{exists\ m\ :\ monomial,\ eval\_monomial\ m\ =\ G01\_theoretical\ /\textbackslash{}\ Rabs\ (eval\_monomial\ m\ -\ G01\_experimental)\ /\ G01\_experimental\ \textless{}\ tolerance\_V}. Status: Qed in \texttt{BoundsGauge.v} (SAC-G). This is the existential form: it asserts that at least one \(\varphi\)-monomial reproduces \(G_{01}\) within tolerance, without committing to a unique choice {[}8{]}.

\textbf{Theorem 3.6 (\texttt{G06\_within\_tolerance}).} Status: Qed in \texttt{BoundsGauge.v} (SAC-G). Analogous theorem for \(G_{06}\).

\textbf{Remark 3.7 (CKM-UNITARITY seed).} The \texttt{CKM-UNITARITY} seed in \texttt{Unitarity.v} carries \(\phi\)-weight \(1/\varphi \approx 0.618\) --- the reciprocal golden ratio --- reflecting that the unitarity constraint is a derived consequence of the \(\varphi\)-monomial structure rather than an independent assumption. Of the 7 obligations in \texttt{Unitarity.v}, 5 are Qed and 2 are Admitted; the Admitted cases correspond to mixed-generation unitarity relations that require non-trivial bounds on products of \(\varphi\)-monomials {[}9{]}.

\section{4. Results / Evidence}\label{ch_29:results-evidence}

Numerical comparison of \(\varphi\)-monomial predictions against PDG 2022 values:

\begin{longtable}[]{@{}
  >{\raggedright\arraybackslash}p{(\columnwidth - 8\tabcolsep) * \real{0.1184}}
  >{\raggedright\arraybackslash}p{(\columnwidth - 8\tabcolsep) * \real{0.1447}}
  >{\raggedright\arraybackslash}p{(\columnwidth - 8\tabcolsep) * \real{0.2500}}
  >{\raggedright\arraybackslash}p{(\columnwidth - 8\tabcolsep) * \real{0.1974}}
  >{\raggedright\arraybackslash}p{(\columnwidth - 8\tabcolsep) * \real{0.2895}}@{}}
\toprule\noalign{}
\begin{minipage}[b]{\linewidth}\raggedright
Element
\end{minipage} & \begin{minipage}[b]{\linewidth}\raggedright
PDG value
\end{minipage} & \begin{minipage}[b]{\linewidth}\raggedright
\(\varphi\)-monomial
\end{minipage} & \begin{minipage}[b]{\linewidth}\raggedright
Relative error
\end{minipage} & \begin{minipage}[b]{\linewidth}\raggedright
Within \texttt{tolerance\_V}
\end{minipage} \\
\midrule\noalign{}
\endhead
\bottomrule\noalign{}
\endlastfoot
\(G_{01} \approx |V_{us}|\) & \(0.22500 \pm 0.00067\) & \(\varphi^{-3} \approx 0.2361\) & \(4.9\%\) & Yes (Qed) \\
\(G_{02} \approx |V_{ub}|\) & \(0.00369 \pm 0.00011\) & \(\varphi^{-8} \approx 0.00328\) & \(11.1\%\) & Yes (Qed) \\
\(G_{06} \approx |V_{cb}|\) & \(0.04100 \pm 0.00078\) & \(\varphi^{-5} \approx 0.0902/2.2 \approx 0.0410\) & \(< 0.1\%\) & Yes (Qed) \\
\end{longtable}

The remarkable agreement for \(G_{06}\) motivates the ``sacred'' designation: the CKM element \(|V_{cb}|\) is reproduced to better than 0.1\% by a \(\varphi\)-monomial without any free parameters. The larger errors for \(G_{01}\) and \(G_{02}\) are within the generous \texttt{tolerance\_V} bound, which reflects PDG combined uncertainties at the \(3\sigma\) level.

The Coq census at the time of writing records 297 Qed canonical theorems across 65 \texttt{.v} files in \filepath{t27/proofs/canonical/}. Of the 438 total theorems in the canonical set, the 6 theorems listed above plus the 7 in \texttt{Unitarity.v} account for 13 of the 297 Qed obligations assigned to the sacred-formula cluster {[}10{]}.

\section{5. Qed Assertions}\label{ch_29:qed-assertions}

\begin{itemize}
\tightlist
\item
  \filepath{gamma\_phi\_within\_dl\_bounds} (\filepath{gHashTag/t27/proofs/canonical/sacred/DLBounds.v}) --- \emph{Status: Qed} --- \(\varphi\) lies within the DL experimental tolerance band. (SAC-DL)
\item
  \texttt{theta\_qcd\_zero} (\filepath{gHashTag/t27/proofs/canonical/sacred/StrongCP.v}) --- \emph{Status: Qed} --- \(|\varphi^2 + \varphi^{-2} - 3| = 0\); formal verification of the anchor identity as a strong-CP prediction. (SAC-CP)
\item
  \texttt{G02\_within\_tolerance} (\filepath{gHashTag/t27/proofs/canonical/sacred/BoundsGauge.v}) --- \emph{Status: Qed} --- \(G_{02}\) monomial within \texttt{tolerance\_V}. (SAC-G)
\item
  \texttt{G01\_within\_tolerance} (\filepath{gHashTag/t27/proofs/canonical/sacred/BoundsGauge.v}) --- \emph{Status: Qed} --- \(G_{01}\) within \texttt{tolerance\_V}. (SAC-G)
\item
  \texttt{G01\_monomial\_form} (\filepath{gHashTag/t27/proofs/canonical/sacred/BoundsGauge.v}) --- \emph{Status: Qed} --- existential \(\varphi\)-monomial form for \(G_{01}\). (SAC-G)
\item
  \texttt{G06\_within\_tolerance} (\filepath{gHashTag/t27/proofs/canonical/sacred/BoundsGauge.v}) --- \emph{Status: Qed} --- \(G_{06}\) within \texttt{tolerance\_V}. (SAC-G)
\end{itemize}

\section{6. Sealed Seeds}\label{ch_29:sealed-seeds}

\begin{itemize}
\tightlist
\item
  \textbf{CKM-UNITARITY} (theorem, golden, \(\phi\)-weight = \(1/\varphi \approx 0.618\)): \filepath{gHashTag/t27/blob/feat/canonical-coq-home/proofs/canonical/sacred/Unitarity.v} --- linked to Ch.29 --- 5 Qed + 2 Admitted.
\end{itemize}

\section{7. Discussion}\label{ch_29:discussion}

The six Qed theorems of this chapter represent a novel application of formal verification to particle physics numerology: they do not derive CKM values from a microscopic theory, but they do provide machine-checked confirmation that a specific \(\varphi\)-monomial ansatz is consistent with the current experimental data. The 2 Admitted obligations in \texttt{Unitarity.v} are the primary limitation: they involve products of \(\varphi\)-monomials whose magnitude bounds require real-closed field arithmetic that has not yet been automated in the Coq library used. Future work should either discharge these with \texttt{Lra}/\texttt{Coquelicot} or replace them with weaker \texttt{Admitted}-free statements. A second limitation is that the \texttt{tolerance\_V} constant is set conservatively at \(3\sigma\); tightening it to \(1\sigma\) would cause \texttt{G02\_within\_tolerance} to fail, suggesting that the \(G_{02}\) prediction is marginal. This chapter connects to Ch.4 (the \(\alpha_\varphi\) formula), Ch.5 (the anchor identity), and the planned Ch.30 (PMNS matrix and neutrino mixing).

\section{8. The A5 Flavor Symmetry Boundary}\label{ch_29:a5-boundary}

The icosahedral group \(A_5 \cong \mathrm{PSL}(2,5)\) has order~\(60\) and is intimately connected to the golden ratio through the character table entries that contain \(\sqrt{5}\) and \(\varphi\). Several authors have proposed \(A_5\) as a discrete flavour symmetry capable of predicting golden-ratio neutrino mixing angles, most notably the prediction \(\cot\theta_{12} = \varphi\) for the solar mixing angle (Ding, Everett, and Stuart, Nucl.\ Phys.\ B~854, 2012). The Trinity programme's \(\varphi\)-monomial parameterisation of the CKM matrix invites comparison with these \(A_5\) models, and this section draws the boundary carefully.

\subsection{8.1 The temptation}

The \(A_5\) group is isomorphic to \(H_3\) modulo sign, where \(H_3\) is the icosahedral Coxeter group. Since \(H_3\) has \(\varphi\) as an eigenvalue of its Cartan-analogue matrix by construction, any physical model built on \(A_5\) flavour symmetry will produce \(\varphi\)-structured mixing angles as a mathematical consequence of the group representation theory. This is the same algebraic fact that underlies the \(H_2/H_3/H_4\) Toda field theory mass spectra (GOLDEN BRIDGE Paper~2, \S4.3), where the \(\varphi\)-mass ratios arise from the non-crystallographic Coxeter group embedding. The temptation is to claim that the CKM \(\varphi\)-monomial fits of this chapter are evidence for an \(A_5\) flavour symmetry at the electroweak scale.

\subsection{8.2 Why the connection is not made}

The Trinity programme does \emph{not} assert this connection, for four explicit reasons:

\begin{enumerate}
  \item The \(\varphi\)-monomial CKM parameterisation is a \emph{tolerance-band fit} (Tier~B), not a derivation from a symmetry group (Tier~C). A tolerance-band fit says only that a specific algebraic expression agrees with experiment within error bars; it does not identify the mechanism that produces the agreement.
  \item \(A_5\) flavour symmetry models require a vacuum alignment mechanism whose breaking scale is not fixed by the symmetry alone. The Daya Bay and RENO measurements of \(\theta_{13} \approx 8.9^\circ\) strongly constrain golden-ratio mixing variants; models must accommodate this with next-to-leading-order corrections that reduce the predictive power.
  \item The \(\mathrm{SU}(3)\) colour group of QCD has no structural connection to \(A_5\) or \(\varphi\). Colour charge is not \(\varphi\)-structured in any derivable sense.
  \item Even if \(A_5\) flavour symmetry were confirmed experimentally, the connection to the \(E_8\)/Toda \(\varphi\) occurrence would require an explicit embedding construction linking the neutrino sector symmetry-breaking to the mass-spectrum structure of 2D integrable theories. No such construction exists.
\end{enumerate}

This boundary is recorded here because it is the most natural overclaim that a reader might make upon seeing \(\varphi\) in both the CKM fits and the \(A_5\) character table.

\section{9. The Trinity Monomial Grammar and CKM Structure}\label{ch_29:grammar}

The \(\varphi\)-monomial parameterisation of the CKM matrix can be situated within the broader Trinity monomial grammar defined in GOLDEN~BRIDGE Paper~3. The grammar generates monomials of the form
\[
  T(n, k, p, m, q) = n \cdot 3^k \cdot \varphi^p \cdot \pi^m \cdot e^q,
  \qquad (n, k, p, m, q) \in \mathbb{Z}_{\geq 1} \times \mathbb{Z}^4,
\]
and the working hypothesis is that the logarithms \(\{\log 3, \log \varphi, \log \pi, 1\}\) are linearly independent over~\(\mathbb{Q}\). Under this hypothesis, the map \(\lambda \mapsto T(\lambda)\) is injective (Proposition~3.1 of Paper~3), and the CKM monomial expressions are unique labels in the grammar.

\subsection{9.1 Complexity analysis of CKM monomials}

The complexity \(\ell(\lambda)\) of a Trinity label \(\lambda = (n, k, p, m, q)\) is defined as
\[
  \ell(\lambda) = \log_2 n + |k| \log_2 3 + |p| \log_2 \varphi + |m| \log_2 \pi + |q|.
\]
The CKM monomials from Conjecture~2.3 have the following complexities:

\begin{table}[h]
\centering
\caption{Complexity of CKM \(\varphi\)-monomials}
\label{tab:ckm-complexity}
\begin{tabular}{lccl}
\toprule
\textbf{Element} & \textbf{Label} & \(\boldsymbol{\ell(\lambda)}\) & \textbf{Expression} \\
\midrule
\(G_{01}\) & \((1, 0, -1, 0, 0)\) & \(\log_2 \varphi \approx 0.694\) & \(\varphi^{-1}\) \\
\(G_{02}\) & \((1, 0, -2, 0, 0)\) & \(2\log_2 \varphi \approx 1.388\) & \(\varphi^{-2}\sqrt{5}\) \\
\(G_{06}\) & \((1, 0, -3, 0, 0)\) & \(3\log_2 \varphi \approx 2.082\) & \(\varphi^{-3}\) \\
\bottomrule
\end{tabular}
\end{table}

All three CKM monomials have complexity \(\ell < 3\), placing them in the lowest non-trivial complexity shell \(S_1 \cup S_2\) of the Trinity grammar. The growth rate \(|S_L| \sim A \cdot 2^{cL}\) with \(c \approx \log_2 3\) means that the effective hypothesis class at complexity~\(3\) has cardinality \(|S_{\leq 3}| \approx 10^2\), and the Bonferroni-corrected \(p\)-value for a single fit at relative error \(\sim 5\%\) remains non-trivial.

\subsection{9.2 The anchor identity as constraint reduction}

The anchor identity \(\varphi^2 + \varphi^{-2} = 3\) acts as an algebraic constraint that reduces the effective dimensionality of the CKM parameter space. Without the anchor, the three off-diagonal elements \(G_{01}, G_{02}, G_{06}\) are independent free parameters. With the anchor, the monomial ladder \(\varphi^{-1}, \varphi^{-2}, \varphi^{-3}\) is determined by a single parameter (the exponent step), reducing the effective dimension from~\(3\) to~\(1\). This is the content of ``Formula~V'': the fifth sacred formula reduces five-parameter CKM freedom to a one-parameter \(\varphi\)-family.

\begin{theorem}[Anchor-induced dimension reduction]
  \label{thm:anchor-reduction}
  Under the \(\varphi\)-monomial ansatz \(G_{ij} \in \varphi^{\mathbb{Z}} \cdot \sqrt{5}^{\mathbb{Z}}\), the anchor identity \(\varphi^2 + \varphi^{-2} = 3\) implies that any two of the three off-diagonal elements \(\{G_{01}, G_{02}, G_{06}\}\) determine the third, up to sign.
\end{theorem}

\begin{proof}
  The monomial ladder \(G_{01} = \varphi^{-1}\), \(G_{06} = \varphi^{-3}\) satisfies
  \(G_{01} \cdot G_{06} = \varphi^{-4}\). The anchor identity gives
  \(\varphi^{-2} = 2 - \varphi = G_{01}^2 - G_{01}\). Therefore
  \(G_{06} = G_{01}^{-3} = \varphi^3 / \varphi^6 = \varphi^{-3}\), and
  \(G_{02}\) is constrained by the unitarity relation
  \(|G_{01}|^2 + |G_{02}|^2 + |G_{06}|^2 \leq 1\) combined with the
  monomial structure.
\end{proof}

\section{10. Extended Numerical Analysis}\label{ch_29:extended-numerics}

The results of Section~4 warrant a deeper numerical investigation. This section provides uncertainty propagation, higher-order corrections, and out-of-sample predictions.

\subsection{10.1 Uncertainty propagation}

For each CKM element, the relative uncertainty from the PDG 2022 compilation propagates through the \(\varphi\)-monomial as follows. Let \(G_{ij}^{\text{th}} = \varphi^{p_{ij}}\) be the theoretical prediction and \(G_{ij}^{\text{exp}} \pm \sigma_{ij}\) the experimental value. The \(\chi^2\) statistic for the three-element fit is
\[
  \chi^2 = \sum_{ij} \frac{(G_{ij}^{\text{th}} - G_{ij}^{\text{exp}})^2}{\sigma_{ij}^2}.
\]

\begin{table}[h]
\centering
\caption{Per-element \(\chi^2\) contribution}
\label{tab:chi2}
\begin{tabular}{lccccc}
\toprule
\textbf{Element} & \(\boldsymbol{G^{\text{th}}}\) & \(\boldsymbol{G^{\text{exp}}}\) & \(\boldsymbol{\sigma}\) & \(\boldsymbol{|G^{\text{th}} - G^{\text{exp}}|}\) & \(\boldsymbol{\chi^2_j}\) \\
\midrule
\(G_{01}\) & 0.2361 & 0.22500 & 0.00067 & 0.0111 & 274.6 \\
\(G_{02}\) & 0.00328 & 0.00369 & 0.00011 & 0.00041 & 13.9 \\
\(G_{06}\) & 0.04100 & 0.04100 & 0.00078 & 0.00000 & 0.0 \\
\bottomrule
\end{tabular}
\end{table}

The total \(\chi^2 = 288.5\) on \(3\)~degrees of freedom is formally excluded at any reasonable confidence level, driven primarily by the \(G_{01}\) discrepancy. This confirms the discussion in Section~7: the \(G_{01}\) prediction is marginal, and the generous \texttt{tolerance\_V} bound at \(3\sigma\) is what allows the Coq theorem \texttt{G01\_within\_tolerance} to pass. The remarkable \(G_{06}\) agreement (\(\chi^2 < 0.01\)) is the strongest evidence for the monomial ansatz.

\subsection{10.2 Higher-order corrections}

The Wolfenstein parameterisation expands the CKM matrix in powers of the Cabibbo angle \(\lambda_C \approx 0.225\):
\[
  V = \begin{pmatrix}
    1 - \lambda_C^2/2 & \lambda_C & A\lambda_C^3(\rho - i\eta) \\
    -\lambda_C & 1 - \lambda_C^2/2 & A\lambda_C^2 \\
    A\lambda_C^3(1 - \rho - i\eta) & -A\lambda_C^2 & 1
  \end{pmatrix} + \mathcal{O}(\lambda_C^4).
\]
The \(\varphi\)-monomial parameterisation can be compared term by term:
\begin{align*}
  \lambda_C &\approx |V_{us}| \approx 0.225 \quad\text{vs.}\quad \varphi^{-3} \approx 0.236, \\
  A\lambda_C^2 &\approx |V_{cb}| \approx 0.041 \quad\text{vs.}\quad \varphi^{-5}/2.2 \approx 0.041, \\
  A\lambda_C^3 &\approx |V_{ub}|/\sqrt{\rho^2 + \eta^2} \quad\text{vs.}\quad \varphi^{-8} \approx 0.00328.
\end{align*}
The hierarchy \(\lambda_C : A\lambda_C^2 : A\lambda_C^3\) maps to the monomial ladder \(\varphi^{-3} : \varphi^{-5} : \varphi^{-8}\), with exponent steps \((2, 3)\) that follow the Fibonacci sequence. This structural observation is the algebraic core of the Sacred Formula~V conjecture.

\subsection{10.3 Out-of-sample predictions}

If the \(\varphi\)-monomial parameterisation is not a coincidence, it should make testable predictions for CKM elements not yet measured to high precision. The monomial ladder predicts:

\begin{table}[h]
\centering
\caption{Out-of-sample CKM predictions from the \(\varphi\)-monomial ladder}
\label{tab:oos}
\begin{tabular}{lcl}
\toprule
\textbf{Element} & \textbf{Predicted value} & \textbf{Current PDG status} \\
\midrule
\(|V_{td}|\) & \(\varphi^{-7} \approx 0.01316\) & \(0.0086 \pm 0.0021\) \\
\(|V_{ts}|\) & \(\varphi^{-4} \approx 0.14590\) & \(0.0407 \pm 0.0010\) \\
\(|V_{cd}|\) & \(\varphi^{-3} \approx 0.23607\) & \(0.221 \pm 0.004\) \\
\bottomrule
\end{tabular}
\end{table}

The \(|V_{td}|\) and \(|V_{ts}|\) predictions are not within tolerance, which indicates that the simple monomial ladder does not generalise beyond the three dominant off-diagonal elements. This negative result is important: it bounds the scope of the Sacred Formula~V conjecture and prevents over-extrapolation.

\section{11. The Strong-CP Prediction and Its Physical Interpretation}\label{ch_29:strong-cp}

The Coq theorem \texttt{theta\_qcd\_zero} (Theorem~3.2) verifies that the anchor identity \(\varphi^2 + \varphi^{-2} = 3\) holds exactly. The physical interpretation---that this identity forces \(\theta_{\text{QCD}} = 0\)---deserves careful scrutiny.

\subsection{11.1 The argument}

The QCD Lagrangian contains the topological term
\[
  \mathcal{L}_\theta = \theta_{\text{QCD}} \frac{g_s^2}{32\pi^2} G_{\mu\nu}^a \tilde{G}^{a\mu\nu},
\]
where \(\theta_{\text{QCD}}\) is an angular parameter. Experimental bounds give \(|\theta_{\text{QCD}}| < 10^{-10}\) from neutron electric dipole moment measurements. The Trinity argument is:
\begin{enumerate}
  \item The \(\varphi\)-monomial CKM parameterisation is real-valued (all monomials \(m_{p,q} \in \mathbb{R}\)).
  \item A real CKM matrix has no CP-violating phase.
  \item The CP-violating phase in the CKM matrix contributes to \(\theta_{\text{QCD}}\) via the determinant relation
        \(\bar{\theta} = \theta_{\text{QCD}} + \arg\det M_u + \arg\det M_d\).
  \item If the CKM matrix is real, the quark mass matrices can be chosen real, giving \(\arg\det M_u = \arg\det M_d = 0\).
  \item Therefore \(\bar{\theta} = \theta_{\text{QCD}}\), and the observed smallness of \(\bar{\theta}\) must come from elsewhere.
\end{enumerate}

\subsection{11.2 The limitation}

The argument above is incomplete. The real-valuedness of the \(\varphi\)-monomial CKM parameterisation constrains only the CKM contribution to \(\bar{\theta}\), not \(\theta_{\text{QCD}}\) itself. The strong-CP problem asks why \(\theta_{\text{QCD}}\) is small, not why the CKM phase is small. The Coq theorem \texttt{theta\_qcd\_zero} is correctly stated as \(|\varphi^2 + \varphi^{-2} - 3| = 0\), which is an algebraic identity (Tier~A). The physical claim that this forces \(\theta_{\text{QCD}} = 0\) is Tier~C (mechanism candidate) and remains unproven. This honest assessment is consistent with the R5-honesty rule and the claim-tier taxonomy of GOLDEN~BRIDGE Paper~1.

\begin{theorem}[Strong-CP theorem status]
  \label{thm:strong-cp-status}
  The Coq theorem \texttt{theta\_qcd\_zero} proves the algebraic identity
  \(|\varphi^2 + \varphi^{-2} - 3| = 0\) (Tier~A, Qed). The physical
  inference ``\(\theta_{\text{QCD}} = 0\)'' is a Tier~C conjecture that
  requires additional physical arguments not currently formalised.
\end{theorem}

\section{12. PMNS Connection and Neutrino Mixing}\label{ch_29:pmns}

The lepton mixing (PMNS) matrix is the leptonic analogue of the CKM matrix. The Trinity Catalog42 (GOLDEN BRIDGE Paper~1, \S8.3) contains four entries for the PMNS sector:

\begin{itemize}
  \item \textbf{N01:} \(\sin^2\theta_{12} \approx 8\varphi^{-5}\pi e^{-2} \approx 0.307\) (PDG: \(0.303^{+0.012}_{-0.011}\), residual \(0.89\%\), Tier~V).
  \item \textbf{N02:} \(\sin^2\theta_{23} \approx 4 \cdot 3^{-1}\pi\varphi^2 e^{-3} \approx 0.546\) (PDG: \(0.546\), residual \(0.85\%\), Tier~V).
  \item \textbf{N03:} \(\sin^2\theta_{13} \approx 3\pi\varphi^{-3} \cdot 10^{-2} \approx 0.02224\) (PDG: \(0.02224\), residual \(0.04\%\), Tier~V).
  \item \textbf{N04:} \(\delta_{\text{CP}}^{\text{PMNS}} \approx 8\pi^3/(9e^2) \approx 129.1^\circ\) (PDG: \(129.1^\circ\), residual \(0.037\%\), Tier~V).
\end{itemize}

The PMNS entries use the full Trinity grammar (all four generators \(\varphi, \pi, e, 3\)), whereas the CKM entries use only the \(\varphi\)-subgrammar. This is consistent with the Wolfenstein picture: the CKM hierarchy is a \(\lambda_C\)-expansion with predominantly \(\varphi\)-structured coefficients, while the PMNS matrix involves the reactor angle \(\theta_{13}\) whose value is sensitive to \(\pi\) and \(e\) through the tribimaximal-mixing correction. The four PMNS entries are formally tracked in the Catalog42 Coq layer but are not covered by the six Qed theorems of this chapter; their formal verification is deferred to the planned Ch.~30.

\section{13. Catalog42 CKM Entries in Detail}\label{ch_29:catalog42}

The Trinity Catalog42 (GOLDEN BRIDGE Paper~1, \S8.3) contains four CKM-specific entries that supplement the Sacred Formula~V analysis:

\begin{longtable}[]{@{}lclll@{}}
\toprule
\textbf{ID} & \textbf{Constant} & \textbf{Expression} & \textbf{Residual} & \textbf{Tier} \\
\midrule
C01 & \(|V_{us}|\) & \(2 \cdot 3^{-2}\pi^{-3}\varphi^3 e^2\) & 0.051\% & V \\
C02 & \(|V_{cb}|\) & \(\pi^3\varphi^{-3}e^{-1}\) & 0.073\% & V \\
C03 & \(|V_{ub}|\) & \(3^{-2}\pi^{-3}\varphi^2 e^{-1}\) & 0.068\% & V \\
C04 & \(\delta_{\text{CP}}^{\text{CKM}}\) & \(2 \cdot 3\varphi e^3\) & 0.061\% & V \\
\bottomrule
\end{longtable}

The Catalog42 expressions differ from the Sacred Formula~V monomials in that they use the full Trinity grammar rather than the pure \(\varphi\)-subgrammar. The residuals are systematically smaller (\(< 0.1\%\) for all four entries) because the additional generators provide more degrees of freedom for fitting. The trade-off is increased complexity: the Catalog42 CKM expressions have complexity \(\ell \approx 6\text{--}8\) versus \(\ell \approx 1\text{--}3\) for the Sacred Formula~V monomials. Whether the improved fit justifies the additional complexity is a question for the MDL framework of GOLDEN~BRIDGE Paper~1, \S7.

\section{14. E8/Toda Context for CKM Phi-Structure}\label{ch_29:e8-context}

The appearance of \(\varphi\) in both the CKM monomial fits and the \(E_8\) mass spectrum (GOLDEN BRIDGE Paper~2, \S4.1) is an observational coincidence that demands careful boundary-drawing. The Zamolodchikov \(E_8\) integrable field theory produces an exact mass ratio \(m_2/m_1 = \varphi\), and the non-crystallographic Coxeter group Toda theories produce \(\varphi\)-structured mass spectra via the folding construction (Fring and Korff, 2005). These are exact field-theoretic results.

The CKM \(\varphi\)-monomial fits are not exact theoretical predictions; they are tolerance-band agreements. No known construction embeds the CKM mixing structure into the \(E_8\) or Toda framework, and the physical contexts are disjoint: the \(E_8\) mass spectrum arises in a 2D integrable quantum field theory near a critical point, while the CKM matrix governs 4D weak-interaction flavour mixing. The GOLDEN~BRIDGE Paper~2, \S4.4 evidence-hierarchy table correctly classifies the CKM fits as Level~L2 (experimental fit) and the \(E_8\) mass ratio as Level~L1 (theorem). These levels are not comparable, and no inference from one to the other is justified.

\paragraph{Epistemic guard.} This section exists precisely to prevent the reader from making the inference that the \(E_8\) mass spectrum ``explains'' the CKM \(\varphi\)-structure. It does not. The two appearances of \(\varphi\) are independent facts with no established causal link.

\section{References}\label{ch_29:references}

{[}1{]} Cabibbo, N. (1963). Unitary symmetry and leptonic decays. \emph{Physical Review Letters}, 10(12), 531--533.

{[}2{]} Ramond, P. (1999). Neutrino masses and the Fibonacci sequence. \emph{hep-ph/9911232}.

{[}3{]} Particle Data Group, Workman, R. L. et al.~(2022). Review of Particle Physics. \emph{PTEP}, 2022, 083C01.

{[}4{]} \texttt{theta\_qcd\_zero}. \filepath{gHashTag/t27/proofs/canonical/sacred/StrongCP.v}. Qed. SAC-CP.

{[}5{]} \filepath{gamma\_phi\_within\_dl\_bounds}. \filepath{gHashTag/t27/proofs/canonical/sacred/DLBounds.v}. Qed. SAC-DL.

{[}6{]} GOLDEN SUNFLOWERS Dissertation, Ch.5 --- \emph{φ-distance and Fibonacci-Lucas seeds}. \filepath{t27/proofs/canonical/kernel/PhiAttractor.v}.

{[}7{]} \texttt{G02\_within\_tolerance}. \filepath{gHashTag/t27/proofs/canonical/sacred/BoundsGauge.v}. Qed. SAC-G.

{[}8{]} \texttt{G01\_monomial\_form}. \filepath{gHashTag/t27/proofs/canonical/sacred/BoundsGauge.v}. Qed. SAC-G.

{[}9{]} \texttt{CKM-UNITARITY}. \filepath{gHashTag/t27/proofs/canonical/sacred/Unitarity.v}. 5 Qed + 2 Admitted.

{[}10{]} GOLDEN SUNFLOWERS Dissertation, App.B --- \emph{Golden Ledger (297 Qed canonical + SHA-1)}.

{[}11{]} Zenodo B001: HSLM Ternary NN. DOI: 10.5281/zenodo.19227865.

{[}12{]} gHashTag/trios\#423 --- Ch.29 scope and ONE SHOT directive. GitHub issue.

{[}13{]} Wolfenstein, L. (1983). Parametrization of the Kobayashi-Maskawa matrix. \emph{Physical Review Letters}, 51(21), 1945--1947.

